\section{Huffman 编码与贪心最优性}
\label{sec:07-Information-Theory/04-Source-Coding/02}

上一节把``一套码''和``一个概率分布''对上了号，理想码长是 \(-\log_2 p_i\)。可理想码长几乎从不是整数，而实际码字的长度只能是整数。于是问题变成：在 Kraft 不等式圈出的那些整数长度组合里，哪一组让平均码长最小？

先看看这个问题为什么不能逐个符号独立回答。你或许想给每个符号单独挑一个最接近 \(-\log_2 p_i\) 的整数，可码长不是各自为政的量——它们共用一棵树。把某片叶子上移一层，它腾出的位置和封锁的子树都会变，其他符号的可选深度跟着改。局部看起来的最优选择，放回树里往往互相打架。

Huffman 的做法是把这个全局的树形优化，压成一串只看两个数的局部决策，而且这个压缩是精确的，不是启发式。

\subsection{最优树的形状先被两条约束框死}

盯住任意一棵最优的二元前缀树。若某个内部节点只有一个孩子，把那棵子树整体上提一层，所有相关码字都变短，没有码字变长——这与最优矛盾。所以最优树必是满树，每个内部节点恰有两个孩子。

满树推出第二条：最深那一层的叶子必然成对出现。最深叶的同胞不可能是内部节点，否则那个内部节点底下还得再长出更深的叶子，与``最深''冲突。

于是最深处有一对同胞叶。它们该分给谁？答案是概率最小的那两个符号，理由是一个交换动作。假设最深的一对同胞上坐着符号 \(x\)，而某个更浅的位置坐着概率更小的符号 \(a\)。把两者对调，平均码长的变化是

\[
(p_a-p_x)\bigl(\ell_{\text{深}}-\ell_{\text{浅}}\bigr)\le 0,
\]

两个因子一负一正，乘积非正。交换只会让平均码长不增，所以总存在一棵最优树，把最小的两个概率放在最深的同胞位置上。这类论证不构造新解，只是把已知的最优解重新排布成自己想要的样子；Kruskal 的最小生成树和不少调度问题的最优性证明用的是同一招。

\subsection{合并一次，问题就小一格}

既然最小的两个符号 \(a,b\) 可以安置成最深的同胞，就可以把它们捏成一个概率为 \(p_a+p_b\) 的复合符号。原来给 \(m\) 个符号建最优树的问题，缩成给 \(m-1\) 个符号建最优树。

两个问题的代价关系是干净的。设复合符号在缩约树里的码长为 \(r\)，拆开后两个子符号的码长都是 \(r+1\)，于是

\[
L=L'+(p_a+p_b).
\]

多出来的那一项与缩约树的具体形状无关。这正是贪心成立的关键：无论下层怎么排，合并这一步付的代价都一样，所以让缩约树最优就让原树最优。反过来说，若原问题有更好的树，把它调整成最小两个符号相邻的形式再合并，就会得到一棵比 \(L'\) 更好的缩约树，与缩约树的最优性矛盾。对符号数做归纳，基例是两个符号——只有一种满树，各占一个比特——整个算法的最优性就落地了。

算法本身因此只剩一句话：反复取出当前最小的两个权重，合并成一个新节点放回去，直到只剩树根；再从根往下给左右分支标 \(0\) 和 \(1\)，读出每片叶子的码字。

拿概率 \(0.4,0.3,0.2,0.1\) 的四个符号试一次，得到的一种码是 \(A\mapsto0\)、\(B\mapsto10\)、\(C\mapsto110\)、\(D\mapsto111\)，平均码长 \(1.9\) 比特，而熵是 \(1.846\) 比特。合并过程中出现过两个相等的权重 \(0.3\)，选哪个先合会长出不同形状的树、给出不同的码字，但平均码长相同。Huffman 保证的是最优值，不是最优解唯一。

\subsection{取整之后，码长与自信息的对照}

把这四个符号的理想码长和 Huffman 给的整数码长并排画出来，取整这件事的形状就清楚了。

\begin{center}
\begin{tikzpicture}[x=1cm,y=0.82cm,font=\small,>=stealth]
  \draw[->,black!55] (0.1,0) -- (8.4,0);
  \draw[->,black!55] (0.1,0) -- (0.1,4.0) node[above,font=\footnotesize] {比特};
  \foreach \y in {1,2,3} {
    \draw[black!20] (0.1,\y) -- (8.2,\y);
    \node[left,font=\footnotesize,black!60] at (0.05,\y) {\y};
  }
  \fill[blue!45] (0.4,0) rectangle (0.95,1.322);
  \fill[black!30] (1.05,0) rectangle (1.6,1.0);
  \fill[blue!45] (2.4,0) rectangle (2.95,1.737);
  \fill[black!30] (3.05,0) rectangle (3.6,2.0);
  \fill[blue!45] (4.4,0) rectangle (4.95,2.322);
  \fill[black!30] (5.05,0) rectangle (5.6,3.0);
  \fill[blue!45] (6.4,0) rectangle (6.95,3.322);
  \fill[black!30] (7.05,0) rectangle (7.6,3.0);
  \node[below,font=\footnotesize] at (1.0,-0.08) {\(A\)  \(0.4\)};
  \node[below,font=\footnotesize] at (3.0,-0.08) {\(B\)  \(0.3\)};
  \node[below,font=\footnotesize] at (5.0,-0.08) {\(C\)  \(0.2\)};
  \node[below,font=\footnotesize] at (7.0,-0.08) {\(D\)  \(0.1\)};
  \fill[blue!45] (0.45,3.7) rectangle (0.75,3.9);
  \node[right,font=\footnotesize] at (0.8,3.8) {\(-\log_2 p_i\)};
  \fill[black!30] (0.45,3.3) rectangle (0.75,3.5);
  \node[right,font=\footnotesize] at (0.8,3.4) {Huffman 码长};
  \draw[red!70,<->] (7.72,3.0) -- (7.72,3.322);
  \node[right,font=\footnotesize,red!70,align=left] at (7.78,3.16) {比自信息\\还短 \(0.32\)};
\end{tikzpicture}

\emph{理想码长一般不是整数，Huffman 只能在可行的整数组合里选。注意最右边的 \(D\)：它拿到的码字比自己的自信息还短。单个码字完全可以短于 \(-\log_2 p_i\)，被熵压住的从来只是加权平均。}
\end{center}

这张图顺带纠正一个常见的误读：熵下界不是对每个码字说的，是对平均说的。低概率符号占的权重小，即使它的码长偏短，对平均的贡献也有限；账要在整个分布上算。

\subsection{离熵有多远，以及什么时候恰好贴上}

Huffman 是前缀码里的最优者，因此下界照收；上界可以借 Shannon 码的存在性——Huffman 不会比它差。两头夹起来就是

\[
H(X)\le L_{\mathrm H}<H(X)+1.
\]

上一节说过取等的条件：所有概率都是 \(2^{-k}\) 的形式。比如 \((1/2,1/4,1/8,1/8)\)，理想长度 \((1,2,3,3)\) 全是整数，Huffman 精确打到 \(1.75\) 比特。除此之外总有零头，而零头来自整数约束这个外部限制，不是算法的缺陷。

差额最难看的时候，是符号少而分布偏。Bernoulli\((0.99)\) 只有两个符号，无论怎么分配都得给每个符号一个比特，而熵只有 \(0.081\) 比特，浪费超过十倍。Huffman 在自己的规则内已经无懈可击，问题在规则本身：逐符号编码强制每个符号至少消耗一整个比特，哪怕它几乎不携带不确定性。

破局有两条路。把 \(k\) 个符号绑成一个超级符号做块编码，取整的浪费被摊到 \(k\) 个符号上，每符号冗余降到 \(1/k\) 以下；代价是码表规模按 \(m^k\) 膨胀，\(256\) 个符号的字母表绑两个就是六万多行，绑三个已不可行。另一条路是不再给每个块分配整数长度，改为把整段序列映到一个区间，让分数比特在序列内部自由累积——下一节走的就是这条。

\subsection{最优性的有效期只覆盖模型内部}

真实压缩器还要回答几个平均码长公式管不到的问题。解码端并不知道概率表，码表本身就得编码传过去；小文件里，码表可能比正文还大。Canonical Huffman 通过约定同长度码字按字母序递增来消除这份开销的大头，只传每个符号的码长序列而不是整棵树。若改用自适应方案边编码边更新概率，码表不用传了，但编解码两端的模型必须逐比特同步——一处偏差之后的所有输出都不可恢复。至于并列权重怎么打破，本身无所谓对错，只要两端用同一条规则。

字母表更大时结论照旧成立。\(D\) 元 Huffman 每步合并最小的 \(D\) 个权重，唯一的额外手续是满 \(D\) 叉树的叶子数必须形如 \((D-1)k+1\)，不够就补几个概率为零的虚拟符号占位，它们不参与传输。

更深一层的限制在于概率表本身。Huffman 对编码时用的那个分布最优；数据的分布若在中途漂移，最优性立刻失效，而算法不会报警，只会安静地变长。这不是对 Huffman 的否定，而是所有``最优''结论的通用条款：保证写在一个精确声明的模型内部，走出模型一步，保证就作废。知道那条线画在哪里，才谈得上信任它。

下一节把``每个符号一个整数长度码字''这条规则整个拆掉。一旦允许零头跨符号累积，取整只在整段消息的末尾发生一次，逐符号编码那个恼人的一比特就变成了整条消息上的两个比特。
